\section{Discussion and limitations}
\label{sec:discussion}

\paragraph{One question, two axes.} Both experiments degrade the same
resource---globally informed, per-step credit assignment---along orthogonal
axes. Axis~1 fixes \emph{when} (every step) and varies \emph{what} each
synapse may know; the answer is that connectivity search survives the loss of
weight transport and the chained backward (\dfa{}), and fails gracefully but
substantially when compressed to a scalar (\rmhebb{}). Axis~2 fixes
\emph{what} (exact gradients) and varies \emph{when} depth is paid for; the
answer is that ${\sim}4\%$ of steps, placed adaptively, recover roughly
two-thirds of the accuracy gap between never and always paying for
depth---and that among the schedules we tested at that budget, only the
adaptive ones do. In both cases the interesting quantity is not the best
achievable accuracy but the \emph{shape of the degradation}: graceful along
locality until the vector-to-scalar cliff, graceful along frequency provided
the schedule is closed-loop.

\paragraph{What we do not claim.} All experiments use MLPs on MNIST and
CIFAR-10 with three seeds; absolute accuracies are far from competitive, by
design. We make no claim that these orderings persist in convolutional or
attention-based architectures, at scale, or under distribution shift. The
per-step timing model transfers to transformers only in outline (a
last-block-plus-head front step has the same structure, but attention
backward and activation memory change the constants). The \rmhebb{} result
in particular bounds only the specific three-factor rule we stabilized; a
better scalar-modulated rule may exist, and our finding should be read as
``the obvious one is insufficient,'' not as an impossibility result.

\paragraph{What we believe transfers.} Three mechanisms, not the numbers:
(i)~the cost asymmetry of nested versus partitioned blocks---depth truncation
is exclusive to the top block, so compute-aware scheduling should treat block
choice and burst timing as separate decisions; (ii)~burst-as-measurement with
backoff---any scheduler that must decide whether expensive credit assignment
is still productive faces the cried-wolf problem, and exponential backoff is
a simple, robust answer; (iii)~the danger of isolated deep steps into a stale
trunk, which links training stability to the \lpft{} feature-distortion
mechanism and predicts that interleaved partial-training schemes need either
burst structure or gradient clipping.

\paragraph{Open directions.} Making the trigger's implicit estimate of
inequality~\eqref{eq:greedy} explicit (e.g., cheap stochastic probes of
$\|\nabla L\|$) would turn \fw{} from a heuristic into an approximate optimal
stopping rule. On Axis~1, the gap between \dfa{} and \ep{} may close with
layer-wise mask budgets or momentum on scores; and the vector-to-scalar cliff
invites a controlled study of intermediate signals (low-rank error sketches,
per-class scalars) to locate the threshold precisely.

\section{Conclusion}
Random networks already contain their solutions; finding them does not
require backpropagation's full machinery---an error broadcast suffices---but
it does require more than the scalar-modulated rule we could stabilize. And
when exact gradients are available, most steps do not need them: a
plateau-triggered controller with backoff spends $4\%$ of the depth budget to
recover about two-thirds of the accuracy gap between never and always paying
for it, and beats every fixed alternative we tested at that budget. Both results argue that
credit assignment is best understood as a budgeted resource with structure in
\emph{where} and \emph{when} it is spent, not a fixed per-step cost.

\paragraph{Reproducibility.} Code, run logs, and all raw result CSVs are
available at \url{https://github.com/XTAM-AI/front-wheeler}. Every number in this paper is the
mean over seeds $\{0,1,2\}$ with the exact commands recorded in the run
scripts.
